\section{Boundary lengths, cone angles, and continuation}
\label{sec:cones}

\subsection{Fixed signature and the polynomial chamber}

Allow $a$ cone points with opening angles $\theta_1,\ldots,\theta_a$ and
$b$ geodesic boundary components with lengths $L_1,\ldots,L_b$, including
zero lengths for cusps. Thus $n=a+b$ and $d=3g-3+n$. The area is
\begin{equation}
 \operatorname{Area}(X)=2\pi(2g-2+n)-\sum_{i=1}^{a}\theta_i>0.
 \label{eq:gb}
\end{equation}
There is no separate puncture parameter in the volume polynomial: a cusp
is represented by $L=0$, or equivalently by cone opening zero.

Use the logarithmic tautological class
$\kappa_1=\pi_*(\psi_{n+1}^2)$ on $\MM_{g,n}$. The fixed-boundary WP class,
continued to cone parameters, is
\begin{equation}
 \Omega_0=2\pi^2\kappa_1+\frac12\sum_{j=1}^{b}L_j^2\psi_j
                          -\frac12\sum_{i=1}^{a}\theta_i^2\psi_i.
 \label{eq:wpclass}
\end{equation}
The labels in the two sums are understood to be disjoint. In its geometric
validity chamber,
\begin{equation}
 V_{g,n}(\lengths,i\angles)=\frac1{d!}
                          \int_{\MM_{g,n}}\Omega_0^d.
 \label{eq:intersectionvolume}
\end{equation}
The boundary formula follows from the symplectic-reduction/intersection
description of WP volumes~\cite{MirzIntersection}; imaginary lengths
$L=i\theta$ account for the minus sign~\cite{DoNorbury}.

For mixed positive boundaries and cone points, a safe geometric regime is
$0\le\theta_i<\pi$. For pure cone/cusp surfaces, the uncorrected formula
holds in the larger chamber where every pair of openings sums to less than
$2\pi$~\cite[Theorem 1.1 and Corollary 1.2]{AnagnostouNorbury}.
In both statements the area condition~\eqref{eq:gb} and the prescribed
labels remain part of the hypotheses. Polynomial evaluation outside these
ranges is always algebraically defined, but requires a separate geometric
justification. In particular, a single cone of arbitrary
opening below $2\pi$, with only cusps as the other ends, does not create a
cone-collision wall.

For example, on $\MM_{0,4}$ one has
$\int\kappa_1=\int\psi_i=1$, so
\begin{equation}
 V_{0,4}(\lengths,i\angles)
 =2\pi^2+\frac12\sum L_j^2-\frac12\sum\theta_i^2
 \label{eq:mixedv04}
\end{equation}
whenever this chamber formula applies. On $\MM_{1,1}$,
$\int\kappa_1=\int\psi_1=1/24$, giving the torus normalization in
\eqref{eq:basevolumes} and its one-cone continuation.

\subsection{Pure cone/cusp volumes beyond the first chamber}

For $0\le\theta_i<2\pi$, put
\[
 w_i=1-\frac{\theta_i}{2\pi},\qquad
 \Phi_S=\sum_{i\in S}\theta_i-2\pi(|S|-1),\qquad u_+=\max(u,0).
\]
Assume $2g-2+\sum_iw_i>0$. A collision of the marked points in $S$ is
allowed in the weighted compactification when $\sum_{i\in S}w_i\le1$.
The class pulled back to ordinary Deligne--Mumford space is
\begin{equation}
 \boxed{
 \Omega_{\angles}=2\pi^2\kappa_1-\frac12\sum_i\theta_i^2\psi_i
      +\frac12\sum_S(\Phi_S)_+^2\,\delta_{0:S}.}
 \label{eq:wallclass}
\end{equation}
The sum runs over stable rational-tail divisors, with $|S|\ge2$. The
geometric volume is $\int\Omega_{\angles}^{d}/d!$. This is the
wall-crossing class of Anagnostou--Mullane--Norbury
\cite[Definitions 2.12--2.13 and Lemma 2.14]{AMN}, also expressed in the
cone-volume formulation of Eberhardt--Turiaci~\cite{ET}.
Here stability means that both components of the nodal curve are stable
in ordinary DM space; in genus zero this requires also $|S^c|\ge2$.
All classes in~\eqref{eq:wallclass}, including $\kappa_1$ and $\psi_i$,
are on that ordinary compactification. They must not be replaced by
weighted-space classes without applying the reduction pullback.

For completeness, the coefficient follows immediately from the weighted
formula: its normalized correction is
$(1-\sum_{i\in S}w_i)^2\delta_{0:S}$, and multiplication by $2\pi^2$
gives $(\Phi_S)^2/2$. If $r$ is the reduction to the relevant weighted
compactification, then $r$ is proper and birational. Therefore
$r_*[\MM_{g,n}]=[\MM_{g,\mathbf w}]$, and the projection formula justifies
integrating the top power of the pulled-back class on ordinary DM space.
On a wall $\Phi_S=0$ the added term vanishes. The formula therefore
extends continuously from adjacent chambers, although its polynomial
expression changes across the wall. In dimension greater than one the
top power includes products and self-intersections of the boundary
divisors; adding their coefficients as scalar volume corrections would
not compute the volume.

In genus zero $\delta_{0:S}=\delta_{0:S^c}$. There is no positive-term
double counting in~\eqref{eq:wallclass}: stability gives $\sum_iw_i>2$,
so both $S$ and its complement cannot have weight at most one.

For four labeled cone/cusp points the formula becomes particularly simple:
\begin{equation}
 \boxed{
 V_{0,4}(i\angles)=2\pi^2-\frac12\sum_{i=1}^4\theta_i^2
 +\frac12\sum_{1\le i<j\le4}(\theta_i+\theta_j-2\pi)_+^2.}
 \label{eq:conev04}
\end{equation}
For openings $(0,0,3\pi/2,3\pi/2)$, the area is positive but the uncorrected
polynomial gives $-\pi^2/4$. The single active correction contributes
$\pi^2/2$, giving the actual positive volume $\pi^2/4$. This example
pinpoints the failure of unrestricted polynomial continuation.

With a positive geodesic boundary and a cone opening greater than $\pi$,
additional degenerations can occur when the cone approaches that boundary.
Anagnostou--Norbury exhibit such a nonnodal degeneration
when the positive boundary length is sufficiently large
\cite[Lemma 2.10]{AnagnostouNorbury}. We consequently assert the full
wall formula here for pure cones/cusps, and the preceding small-cone regime
for arbitrary mixed positive boundaries. An unrestricted mixed-signature
geometric theorem would require its own hypotheses and boundary analysis.

\subsection{What continuation means for compact polyhedra}

The boundary of a compact convex hyperbolic polyhedron has intrinsic cone
opening $\theta_v$ equal to the sum of its face angles at $v$; its singular
curvature is $2\pi-\theta_v$. Neither is an individual dihedral angle.
Thus the substitution $L\mapsto i\theta$ concerns an intrinsic holonomy
parameter. It is not a literal substitution of arbitrary hyperideal vertex
coordinates by compact ones.

At a local algebraic level the continuation is natural. For example,
\begin{equation}
 K_L(s)=\frac{\sinh(L/2-s)}{\sinh(L/2)}
 \quad\longmapsto\quad
 \frac{\sin(\theta/2-\alpha)}{\sin(\theta/2)}
 \quad\text{under }(L,s)=(i\theta,i\alpha).
 \label{eq:kernelcontinuation}
\end{equation}
Trace and Gram formulas have analogous analytic continuations. Recovering
a real geometric volume also requires a real geometric domain, a valid
coordinate chart, and the appropriate collision chamber. The polynomial
example above shows why these cannot be inferred from the formal algebra.

